\chapter{Results}
\label{ch:results}

The results follow the same progression as the tasks themselves. Fixed hover first asks whether
the vehicle can remain calm around one point. Moving-target hover then asks it to travel and
settle repeatedly. Physical landing adds a shrinking time margin, actuator sequencing and
irreversible ground contact. Figure~\ref{fig:result-summary} gives the headline evaluation
outcomes, while the sections that follow show what those aggregates hide.

\begin{figure}[htbp]
  \centering
  \begin{tikzpicture}[node distance=0.38cm]
    \node[rocketbox,text width=0.19\textwidth,minimum height=1.95cm] (hover) {
      \textcolor{rocketTeal}{\LARGE\textbf{250/250}}\\[-1pt]
      \small hover horizons\\[2pt]
      \footnotesize completed for \SI{60}{s}};
    \node[rocketbox,text width=0.19\textwidth,minimum height=1.95cm,right=of hover] (track) {
      \textcolor{rocketTeal}{\LARGE\textbf{76.4\%}}\\[-1pt]
      \small tracking success\\[2pt]
      \footnotesize 191 / 250 trials};
    \node[rocketbox,text width=0.19\textwidth,minimum height=1.95cm,right=of track] (land) {
      \textcolor{rocketBlue}{\LARGE\textbf{84.8\%}}\\[-1pt]
      \small landing success\\[2pt]
      \footnotesize 212 / 250 trials};
    \node[rocketbox,text width=0.19\textwidth,minimum height=1.95cm,right=of land] (feet) {
      \textcolor{rocketAmber}{\LARGE\textbf{36\%}}\\[-1pt]
      \small full-difficulty landing\\[2pt]
      \footnotesize 18 / 50 trials};
  \end{tikzpicture}
  \caption{Evidence summary for the completed deterministic evaluations. The hover count is
  horizon completion, not a success rate, because that task has no binary criterion. Tracking
  and overall landing rates include all five difficulty bands; the last card isolates the
  landing policy's hardest band.}
  \label{fig:result-summary}
\end{figure}

\section{Fixed hover: correcting the vertical failure}

The earlier \texttt{HoverSimple2} run showed why a long, calm-looking episode is not enough: it
stabilized horizontally while slowly climbing away from the requested altitude. The replacement
run, \texttt{HoverFixed}, changed the actuation envelope rather than hiding that result. Minimum
throttle was reduced from 39\% to 25\%, and the unnecessary separate enable action was removed
from the single-engine controller. The derived policy interface consequently changed from four
to three continuous actions and from 30 to 29 observations, so the new run started from fresh
weights.

Figure~\ref{fig:hover-training} follows the resulting 20.08M-step run. Most early episodes still
end in ground impact, but survival then lengthens until the final policy repeatedly reaches the
roughly \SI{300}{s} fuel-depletion boundary. Over the same progression, median vertical error
falls from about \SI{13}{m} to below \SI{1}{m}, vertical speed approaches zero and tilt settles
near \SI{0.3}{\degree}. Horizontal error is less monotonic: it first increases, then falls to
approximately \SIrange{3}{4}{m} near the end. This is the first run in which the physical
quantities jointly describe bounded powered flight near the target altitude.

\begin{figure}[H]
  \centering
  \includegraphics[width=\textwidth]{training/training-hover.pdf}
  \caption{Physical learning curves for \texttt{HoverFixed}. Completed episodes from both
  revisions are ordered chronologically and grouped into 0.5-million-decision-step bins. Lines
  show medians and shading the interquartile range. The final plateau in survival and propellant
  use corresponds to reaching fuel depletion rather than a binary hover-success event.}
  \label{fig:hover-training}
\end{figure}

The complete evaluation confirms the vertical fix under a broader set of initial conditions.
All 250 deterministic trials reached the \SI{60}{s} horizon. They began \SIrange{0.12}{6.61}{m}
from the horizontal target with up to \SI{2.61}{m/s} horizontal speed and
\SI{1.40}{\degree} tilt. Figure~\ref{fig:hover-evaluation} shows the median response and its
interquartile envelope. The commanded descent briefly raises median speed to approximately
\SI{7.5}{m/s}, but vertical error falls from \SI{20}{m} to below \SI{1}{m} and the policy then
damps both translation and attitude. At \SI{60}{s}, median speed is \SI{0.077}{m/s}, median tilt
is \SI{0.105}{\degree} and median vertical error is \SI{0.902}{m}.

\begin{figure}[H]
  \centering
  \includegraphics[width=\textwidth]{training/hover-evaluation.pdf}
  \caption{Time-aligned profiles from all 250 \texttt{HoverFixed} evaluation episodes. Lines are
  medians, shaded areas are interquartile ranges and orange markers show the \SI{60}{s}
  endpoints. The policy corrects the initial \SI{20}{m} altitude error and damps motion, but
  converges to a residual horizontal offset.}
  \label{fig:hover-evaluation}
\end{figure}

The terminal means tell the same story: speed \SI{0.073}{m/s}, horizontal speed
\SI{0.066}{m/s}, vertical speed \SI{0.021}{m/s}, tilt \SI{0.100}{\degree} and angular rate
\SI{0.082}{\degree\per\second}. Mean fuel use was \SI{11724.7}{kg}, and only one episode
recorded an engine restart. The remaining weakness is position rather than stability. Horizontal
error first rises, then converges tightly to a median of \SI{4.76}{m}; only 77 of 250 trials end
closer to the target than they began.

These 250 horizon completions are substantially stronger evidence than the previous four-trial
diagnostic, but they are not a 100\% success rate. The evaluator explicitly records that no
binary hover metric is defined. \texttt{HoverFixed} therefore demonstrates repeatable altitude,
velocity and attitude stabilization with a residual planar offset; declaring a solved hover
task requires an acceptance radius chosen before the next evaluation.

\section{Moving-target hover tracking}

Tracking begins where fixed hover ends. The policy must leave a stable state, travel to a new
point and become stable again before the timeout. This makes raw target distance a misleading
learning curve: the relocation radius expands while the policy improves. Figure~\ref{fig:tracking-training}
therefore pairs curriculum progress and success with direction efficiency and settling quality.
Progress rose to 97.6\%, and the binned training-success rate stayed mostly above 80\% after the
initial learning phase. The policy was learning to travel and settle under increasingly difficult
conditions, although none of these training signals is a held-out probability.

\begin{figure}[H]
  \centering
  \includegraphics[width=\textwidth]{training/training-hover-track.pdf}
  \caption{Training progression for \texttt{HoverTrackSimple3}. Completed episodes are ordered
  chronologically and grouped into 0.5-million-decision-step bins. Lines show the median and
  shading the interquartile range, except that training success is the bin mean. This view avoids
  mixing repeated episode indices from the three trainer revisions.}
  \label{fig:tracking-training}
\end{figure}

The physical quantities in Figure~\ref{fig:tracking-training-physical} supply the missing
context. Episode duration and propellant use rise as the target moves farther away and the
policy must complete more difficult transfers. Mean horizontal error therefore does not fall in
the same way as it did in fixed hover; after the curriculum expands, it stabilizes despite a
larger task envelope. Horizontal speed peaks during the first attempt to move aggressively, then
settles near \SI{1}{m/s}. Read together with the improving success curve, these are signs of a
controller adapting to a growing problem rather than a stationary metric simply getting worse.

\begin{figure}[H]
  \centering
  \includegraphics[width=\textwidth]{training/tracking-training-physical.pdf}
  \caption{Physical cost of the expanding \texttt{HoverTrackSimple3} curriculum. Lines are
  binned medians and shaded areas are interquartile ranges. Longer episodes and greater fuel use
  accompany farther target relocations, while speed and horizontal error settle after the early
  learning transient.}
  \label{fig:tracking-training-physical}
\end{figure}

The completed evaluation produced 191 successful episodes out of 250, or 76.4\% with a 95\%
Wilson interval of 70.8--81.2\%. Table~\ref{tab:tracking-bands} gives the exact counts, while
Figure~\ref{fig:tracking-success} shows how the failure appears in time, captures and fuel. The
aggregate number hides a sharp boundary effect: performance was perfect at difficulties 0 and
0.25, 98\% at 0.5, 84\% at 0.75, and 0\% at 1.0.

\begin{table}[htbp]
  \centering
  \caption{\texttt{HoverTrackSimple3} evaluation by curriculum band. Intervals are 95\% Wilson
  intervals; each band contains 50 episodes.}
  \label{tab:tracking-bands}
  \resizebox{\textwidth}{!}{%
  \begin{tabular}{rrrrrr}
    \toprule
    Difficulty & Successes & Rate & 95\% interval & Mean duration (s) & Mean captures \\
    \midrule
    0.00 & 50/50 & 100\% & 92.9--100\% & 31.60 & 3.00 \\
    0.25 & 50/50 & 100\% & 92.9--100\% & 41.43 & 3.00 \\
    0.50 & 49/50 & 98\%  & 89.5--99.6\% & 53.76 & 2.98 \\
    0.75 & 42/50 & 84\%  & 71.5--91.7\% & 74.06 & 2.76 \\
    1.00 & 0/50  & 0\%   & 0.0--7.1\%   & 35.00 & 0.00 \\
    \midrule
    Overall & 191/250 & 76.4\% & 70.8--81.2\% & 47.17 & 2.35 \\
    \bottomrule
  \end{tabular}
  }
\end{table}

\begin{figure}[htbp]
  \centering
  \includegraphics[width=\textwidth]{training/tracking-evaluation.pdf}
  \caption{Fixed-band tracking evaluation. Panel A shows success with 95\% Wilson intervals and
  Panel B the number of completed captures. In Panels C and D, boxes span the interquartile
  range, centre lines are medians, diamonds are means and whiskers extend to \(1.5\) times the
  interquartile range. The exact \(d=1\) boundary produces no captures and a uniform
  \SI{35}{s} timeout.}
  \label{fig:tracking-success}
\end{figure}

Across all bands, mean fuel use was \SI{8759.76}{kg}. Mean final planar distance was
\SI{5.64}{m}; mean final speed was \SI{0.877}{m/s}, composed of
\SI{-0.047}{m/s} vertical and \SI{0.870}{m/s} horizontal velocity. Mean final tilt was
\SI{0.709}{\degree} and angular rate \SI{0.171}{\degree\per\second}. These terminal means mix
successful captures with timeout states and must not be read as capture thresholds.

The full-difficulty failures all ended by target timeout after \SI{35}{s}, with zero captures.
The policy had reached a recorded nonlinear training progress of 97.6\%, but had not actually
trained at an exact, stable 100\% boundary. The abrupt change suggests a curriculum-boundary or
generalization problem. The evidence does not isolate whether the dominant cause is relocation
radius, tightened thresholds, state distribution, insufficient exposure or their interaction;
claiming a specific cause would require controlled evaluations around $d=1$.

The tracking policy therefore learned the complete travel--settle sequence over most of the
tested curriculum, including 84\% success at $d=0.75$, but not at its exact final boundary. This
was the first clear example of training progress and frozen evaluation telling different stories.
Landing keeps the same need to travel and settle, then adds the requirement to do so before
physical contact.

\section{Physical leg landing}

\subsection{Why L10 had to become L11}

L10 established that PPO could learn navigation, braking and stable four-foot support, but its
solution still depended on rapid engine cycling. As that run was continued from 30M to 79.63M
steps, effective curriculum difficulty rose from 12.3\% to 60.2\% while late-window restart count
rose from 10.1 to 23.7. Its fixed-band evaluation then produced 151 successes in 250 episodes and
none at $d=1$. This combination motivated L11: preserve outcome-based rewards, make continuous
one-engine descent physically feasible, and constrain rapid switching.

L11 started from fresh weights and reached the complete curriculum. Figure~\ref{fig:l11-training-progress}
places its learning history beside L10 instead of showing the final run in isolation. Both runs
first learn basic touchdown, but their actuator strategies then separate: L10's median restart
count climbs as its curriculum becomes harder, whereas L11 spends most of training near zero or
one restart while continuing to reduce contact speed. L11 reaches 100\% effective difficulty and
continues improving after the 90M-step resume. In the final 1,000 completed, non-replay episodes
at exact $d=1$, training success was 90.1\%; the saved runtime state's recent counter was 92.9\%.

\begin{figure}[H]
  \centering
  \includegraphics[width=\textwidth]{training/training-leg-landing.pdf}
  \caption{Chronological comparison of L10 and L11, grouped into
  one-million-decision-step bins. Lines are medians with interquartile bands, except that
  training success is a bin mean. Contact speed uses only episodes that reached touchdown. L10
  ends at 79.63M steps; the dotted line marks the L11 resume at 90M. Manifest times isolate the
  intended runs and incomplete shutdown rows are excluded.}
  \label{fig:l11-training-progress}
\end{figure}

The redesign changed more than the success count. Table~\ref{tab:landing-run-comparison} shows
that L11 increased
overall evaluation success by 24.4 percentage points, produced measurable success at exact full
difficulty, and reduced impact and cycling at approximately the same mean fuel use. The restart
reduction is especially large: 18.46 to 1.56 per episode.

\begin{table}[htbp]
  \centering
  \caption{Completed fixed-band landing evaluations. Contact quantities are means over episodes
  that reached touchdown; both evaluations used seed 20257 and 250 episodes.}
  \label{tab:landing-run-comparison}
  \resizebox{\textwidth}{!}{%
  \begin{tabular}{lrr}
    \toprule
    Metric & L10 baseline & L11 final \\
    \midrule
    Final effective curriculum difficulty & 60.25\% & 100\% \\
    Overall success & 151/250 (60.4\%) & 212/250 (84.8\%) \\
    Full-difficulty success & 0/50 (0\%) & 18/50 (36\%) \\
    Mean first-contact vertical speed & \SI{3.02}{m/s} & \SI{1.33}{m/s} \\
    Mean first-contact horizontal speed & \SI{0.70}{m/s} & \SI{0.279}{m/s} \\
    Mean first-contact tilt & \SI{1.54}{\degree} & \SI{0.637}{\degree} \\
    Mean maximum contact impulse & \SI{97.0}{\kilo\newton\second} & \SI{39.1}{\kilo\newton\second} \\
    Mean engine restarts & 18.46 & 1.56 \\
    Mean fuel used & \SI{4590}{kg} & \SI{4529}{kg} \\
    \bottomrule
  \end{tabular}}
\end{table}

Figure~\ref{fig:landing-evaluation-comparison} shows that this was not merely an easier set of
random trials. L10 and L11 used the same seed and the same 50 cases in each band. L11's advantage
widens with difficulty, while its median restart count, vertical contact speed and impact remain
below L10 in every band. The redesign therefore improved both mission completion and the physical
quality of the solution.

\begin{figure}[H]
  \centering
  \includegraphics[width=\textwidth]{training/landing-evaluation-comparison.pdf}
  \caption{L10--L11 comparison over the same fixed-band evaluation protocol. Success whiskers
  are 95\% Wilson intervals. The other panels show band medians with interquartile whiskers;
  contact quantities include only episodes that reached touchdown. Lines connect band summaries
  for readability and do not imply interpolation between curriculum levels.}
  \label{fig:landing-evaluation-comparison}
\end{figure}

The final ten million steps produced a smaller net quality gain. Against the same
seeded episode set, overall success rose from 82.8\% at 90M to 84.8\% at 100M, success at
$d=0.75$ rose from 90\% to 92\%, and success at $d=1$ rose from 30\% to 36\%. Mean first-contact
vertical speed fell from \SI{1.68}{m/s} to \SI{1.33}{m/s}, successful-landing planar distance
from \SI{0.97}{m} to \SI{0.63}{m}, and mean maximum impulse from
\SI{52.4}{\kilo\newton\second} to \SI{39.1}{\kilo\newton\second}. At $d=1$, six previously
failing seeded cases succeeded, while three previous
successes regressed. The net gain is real for this fixed set, but the 36\% Wilson interval remains
wide at 24.1--49.9\%.

\subsection{Fixed-band evaluation}

The final evaluation completed all 250 episodes and produced 212 stable four-foot touchdowns:
84.8\% with a 95\% Wilson interval of 79.8--88.7\%. Table~\ref{tab:landing-bands} shows perfect
performance in the first two bands, 96\% at $d=0.5$, 92\% at $d=0.75$, and 36\% at $d=1$.
Figure~\ref{fig:landing-evaluation-comparison} makes the improvement over L10 visible across the
entire range and shows that the gain is accompanied by gentler contact and less switching.

\begin{table}[htbp]
  \centering
  \caption{L11 landing evaluation by fixed curriculum band. Each band contains 50 episodes;
  intervals are 95\% Wilson intervals.}
  \label{tab:landing-bands}
  \resizebox{\textwidth}{!}{%
  \begin{tabular}{rrrrllll}
    \toprule
    Difficulty & Successes & Rate & 95\% interval & Hard touchdown & Missed pad & Too far \\
    \midrule
    0.00 & 50/50 & 100\% & 92.9--100\% & 0 & 0 & 0 \\
    0.25 & 50/50 & 100\% & 92.9--100\% & 0 & 0 & 0 \\
    0.50 & 48/50 & 96\% & 86.5--98.9\% & 1 & 0 & 1 \\
    0.75 & 46/50 & 92\% & 81.2--96.8\% & 2 & 2 & 0 \\
    1.00 & 18/50 & 36\% & 24.1--49.9\% & 30 & 1 & 1 \\
    \midrule
    Overall & 212/250 & 84.8\% & 79.8--88.7\% & 33 & 3 & 2 \\
    \bottomrule
  \end{tabular}}
\end{table}

Across all bands, 248 episodes reached touchdown. Their mean first-contact speed was
\SI{1.38}{m/s}, consisting of \SI{1.33}{m/s} vertical and \SI{0.279}{m/s} horizontal speed;
mean contact tilt was \SI{0.637}{\degree} and angular rate
\SI{1.72}{\degree\per\second}. Mean maximum contact impulse was
\SI{39146}{N.s}, mean rebound was approximately \SI{6.5}{\micro\metre}, and mean recorded stable
hold was \SI{0.566}{s}. No episode recorded a structural strike, one recorded a foot outside the
pad, and every successful episode had all four feet supported. Mean fuel use was \SI{4529}{kg}.

\subsection{Remaining failure and deployment gap}

At full difficulty, 30 of the 32 failures were hard touchdowns; every one exceeded the
\SI{1.5}{m/s} vertical-speed limit. In 25 of those 30 cases, vertical speed was the only failed
first-contact condition, and ten were within \SI{0.2}{m/s} of passing. Median touchdown vertical
speed improved from \SI{1.745}{m/s} at 90M to \SI{1.612}{m/s} at 100M. The remaining problem is
therefore usually a narrow final-braking margin rather than the broad navigation failure seen in
L3--L7. Figure~\ref{fig:landing-full-difficulty} shows both that concentration near the vertical
threshold and the two rare, much larger excursions.

\begin{figure}[H]
  \centering
  \includegraphics[width=\textwidth]{training/landing-full-difficulty.pdf}
  \caption{Diagnostics at exact full difficulty. Panel A separates the 50 outcomes per run;
  Panel B gives empirical cumulative distributions of contact vertical speed, zoomed to
  \SIrange{0}{6}{m/s}; Panel C zooms into L11's touchdown threshold region, with two hard
  contacts outside the displayed range; and Panel D pairs the same 50 seeds at 90M and 100M.
  The dashed red lines mark the \SI{1.5}{m/s} vertical limit and the dotted amber line marks the
  \SI{0.75}{m/s} horizontal limit.}
  \label{fig:landing-full-difficulty}
\end{figure}

The larger unresolved issue is the difference between training and deployment. Exact-$d=1$
training success was 89.6\% over the final 5,000 completed non-replay episodes and 90.1\% over the
final 1,000, yet deterministic evaluation succeeded in only 36\%. Incomplete curriculum exposure
cannot explain this gap because L11 reached and remained at full difficulty.

The action interface is a plausible, but not yet proven, cause. PPO explores with stochastic
continuous actions during training, whereas the evaluator requests deterministic inference. The
three engine-enable commands are nevertheless encoded as continuous outputs and converted to
binary state through on/off thresholds of $+0.35$ and $-0.35$. A deterministic mean action can
therefore occupy the dead band or switch differently from sampled training actions. Consistent
with that hypothesis, mean restarts rose from 1.32 to 1.56 overall and from 2.02 to 2.46 at
$d=1$ between 90M and 100M, even as touchdown quality improved. The decisive follow-up is an
A/B evaluation of the same checkpoint and seeds with deterministic inference enabled and
disabled, followed by a discrete or hybrid engine-enable action if the gap persists.

The bounded result is still substantial: L11 controls navigation, descent, final flare and
stable support reliably through $d=0.75$, and it completes more than one third of trials at the
full trained boundary. It is not a flight controller or a solved full-difficulty policy.

\section{Answers to the research questions}

Taken together, the three tasks answer the research questions at different levels. For RQ1, the
snapshot, derived policy schema, immutable revisions, telemetry and evaluator form a complete
path from a configurable Unity experiment to an identifiable result.

For RQ2, PPO produced repeatable fixed-hover stabilization across 250 trials, although a
predeclared success rule and tighter planar position hold are still missing. It completed
moving-target tracking reliably through $d=0.75$ and learned physical four-leg landing across
the complete curriculum. The final landing policy still fell to 36\% success at $d=1$, so
reaching full training difficulty did not make the controller uniformly reliable.

For RQ3, the decisive failures formed the same chain as the simulator's development: reward
accumulation encouraged delay or cheap termination; a spawn-order bug removed the requested
initial velocity; incorrectly balanced terms made crashing cheaper than missing; approximate
aerodynamics produced wrong or excessive forces; and implicit UI and policy contracts made
changes difficult to trace. The architecture is largely the record of how those failures were
made observable and prevented from recurring.
